\chapter{向量}

%% ---- 知识框架 ---- %%
\begin{knowledgeframe}
\begin{itemize}
	\item \textbf{平面向量的线性运算}：加法、减法、数乘，及运算的几何意义（三角形法则、平行四边形法则）。
	\item \textbf{平面向量的数量积}：$\vec{a}\cdot\vec{b}=|\vec{a}||\vec{b}|\cos\theta$，可用于求夹角、判断垂直。
	\item \textbf{向量的坐标表示}：$\vec{a}=(x_1,y_1), \vec{b}=(x_2,y_2)$，则 $\vec{a}\cdot\vec{b}=x_1x_2+y_1y_2$。
	\item \textbf{共线与垂直}：$\vec{a}//\vec{b} \iff x_1y_2=x_2y_1$；$\vec{a}\perp\vec{b} \iff x_1x_2+y_1y_2=0$。
	\item \textbf{空间向量}：共面定理——$\vec{p},\vec{a},\vec{b}$ 共面 $\iff$ 存在实数 $x,y$ 使 $\vec{p}=x\vec{a}+y\vec{b}$（$\vec{a},\vec{b}$ 不共线时）。
	\item \textbf{三角形的心}：外心（到顶点等距）、重心（$\overrightarrow{GA}+\overrightarrow{GB}+\overrightarrow{GC}=\vec{0}$）、垂心（$\overrightarrow{PA}\cdot\overrightarrow{PB}=\overrightarrow{PB}\cdot\overrightarrow{PC}$）、内心。
\end{itemize}
\end{knowledgeframe}

%% ---- 第一节：平面向量的运算 ---- %%
\section{平面向量的运算}

\startexercise

\subsection*{A组}

\begin{choice}[source={2025合肥高三二检}, topic={向量的夹角}]
已知向量 $\vec{e}_1=(1,0)$, $\vec{e}_2=(1,\sqrt{3})$, 设 $\vec{a}=4\vec{e}_1+\vec{e}_2$, $\vec{b}=3\vec{e}_1-\vec{e}_2$, 则 $\vec{a}$ 与 $\vec{b}$ 的夹角为
\begin{tasks}(2)
	\task $\dfrac{\pi}{6}$
	\task $\dfrac{\pi}{4}$
	\task $\dfrac{\pi}{3}$
	\task $\dfrac{2\pi}{3}$
\end{tasks}
\end{choice}
\begin{choice-sol}
\textbf{答案 C}

待补充解析。
\end{choice-sol}

\begin{choice}[source={2025深圳中学高三二轮三阶}, topic={向量集合的交集}]
已知向量集合 $M=\{\vec{a} \mid \vec{a}=(1,2)+\lambda(3,4), \lambda \in \R\}$, $N=\{\vec{a} \mid \vec{a}=(-2,-2)+\lambda(4,5), \lambda \in \R\}$, 则 $M \cap N=$
\begin{tasks}(2)
	\task $\varnothing$
	\task $\{(1,1),(-2,-2)\}$
	\task $\{(1,1)\}$
	\task $\{(-2,-2)\}$
\end{tasks}
\end{choice}
\begin{choice-sol}
待补充解析。
\end{choice-sol}

\begin{choice}[topic={向量叉积与面积}]
在四边形 $ABCD$ 中，若 $\vec{AC} = (1,3), \vec{BD} = (-6,2)$，则该四边形的面积为
\begin{tasks}(4)
	\task $\sqrt{5}$
	\task $2\sqrt{5}$
	\task $10$
	\task $20$
\end{tasks}
\end{choice}
\begin{choice-sol}
\textbf{答案 C}

待补充解析。
\end{choice-sol}


\subsection*{B组}

\begin{choice}[source={2024年北京卷选择题\#5}, topic={向量数量积与条件}]
设 $\vec{a},\vec{b}$ 是向量，``$(\vec{a}+\vec{b})\cdot(\vec{a}-\vec{b})=0$'' 是 ``$\vec{a}=-\vec{b}$ 或 $\vec{a}=\vec{b}$'' 的
\begin{tasks}(2)
	\task 充分不必要条件
	\task 必要不充分条件
	\task 充要条件
	\task 既不充分也不必要条件
\end{tasks}
\end{choice}
\begin{choice-sol}
\textbf{答案 B}

待补充解析。
\end{choice-sol}

\begin{choice}[source={2024年高考全国甲卷理科选择题\#9}, topic={向量的垂直与共线}]
已知向量 $\vec{a}=(x+1,x)$，$\vec{b}=(x,2)$，则
\begin{tasks}(1)
	\task ``$x=-3$'' 是 ``$\vec{a}\perp\vec{b}$'' 的必要条件
	\task ``$x=-3$'' 是 ``$\vec{a}//\vec{b}$'' 的必要条件
	\task ``$x=0$'' 是 ``$\vec{a}\perp\vec{b}$'' 的充分条件
	\task ``$x=-1+\sqrt{3}$'' 是 ``$\vec{a}//\vec{b}$'' 的充分条件
\end{tasks}
\end{choice}
\begin{choice-sol}
\textbf{答案 C}

待补充解析。
\end{choice-sol}

\begin{choice}[topic={直线方向向量与夹角}]
设直线 $l_1$ 的一个方向向量为 $\vec{a}$，直线 $l_2$ 的一个方向向量为 $\vec{b}$。"直线 $l_1$ 与 $l_2$ 所成的角为 $\dfrac{\pi}{3}$" 是 "向量 $\vec{a}$ 与 $\vec{b}$ 的夹角为 $\dfrac{2\pi}{3}$" 的
\begin{tasks}(2)
	\task 充分不必要条件
	\task 必要不充分条件
	\task 充要条件
	\task 既不充分也不必要条件
\end{tasks}
\end{choice}
\begin{choice-sol}
待补充解析。
\end{choice-sol}


\subsubsection*{三角形的心}

\begin{choice}[source={必修二P52页习题}, topic={三角形的外心重心垂心}]
已知点 O、N、P 在 $\triangle ABC$ 所在平面内，且 $\lvert\overrightarrow{OA}\rvert=\lvert\overrightarrow{OB}\rvert=\lvert\overrightarrow{OC}\rvert$，$\overrightarrow{NA} + \overrightarrow{NB} + \overrightarrow{NC} = \vec{0}$，$\overrightarrow{PA} \cdot \overrightarrow{PB} = \overrightarrow{PB} \cdot \overrightarrow{PC} = \overrightarrow{PC} \cdot \overrightarrow{PA}$，则点 O、N、P 依次是 $\triangle ABC$ 的
\begin{tasks}(2)
	\task 重心、外心、垂心
	\task 重心、外心、内心
	\task 外心、重心、垂心
	\task 外心、重心、内心
\end{tasks}
\end{choice}
\begin{choice-sol}
\textbf{答案 C}

待补充解析。
\end{choice-sol}


\subsection*{C组}

\begin{choice}[source={必修二P52页习题}, topic={向量的投影与三角形形状}]
已知非零向量 $\overrightarrow{AB}$ 与 $\overrightarrow{AC}$ 满足 $\dfrac{\overrightarrow{AB} \cdot \overrightarrow{BC}}{\lvert\overrightarrow{AB}\rvert} = \dfrac{\overrightarrow{CA} \cdot \overrightarrow{BC}}{\lvert\overrightarrow{AC}\rvert}$ 且 $\dfrac{\overrightarrow{AB} \cdot \overrightarrow{AC}}{\lvert\overrightarrow{AB}\rvert\lvert\overrightarrow{AC}\rvert} = \dfrac{1}{2}$，则 $\triangle ABC$ 为
\begin{tasks}(2)
	\task 三边均不相等的三角形
	\task 直角三角形
	\task 等腰非等边三角形
	\task 等边三角形
\end{tasks}
\end{choice}
\begin{choice-sol}
\textbf{答案 D}

待补充解析。
\end{choice-sol}

\begin{choice}[topic={空间向量的共面条件}]
设 $\vec{a}, \vec{b}$ 是空间中两个给定的\textbf{非零}向量，$\vec{p}$ 是空间中任意一个向量。"向量 $\vec{p}$ 可以表示为 $\vec{p} = x\vec{a} + y\vec{b}$（其中 $x, y$ 为实数）" 是 "向量 $\vec{p}, \vec{a}, \vec{b}$ 共面" 的
\begin{tasks}(2)
	\task 充分不必要条件
	\task 必要不充分条件
	\task 充要条件
	\task 既不充分也不必要条件
\end{tasks}
\end{choice}
\begin{choice-sol}
\textbf{答案 C}

待补充解析。
\end{choice-sol}


\stopexercise
